% !TEX program = xelatex
% Assignment 6 - GBP-collateralized EUR/USD forward and rate products.
\documentclass[11pt]{article}
\usepackage[a4paper,top=1in,bottom=1in,left=1.1in,right=1.1in]{geometry}
\usepackage{setspace}
\setstretch{1.18}
\setlength{\parindent}{1.5em}
\setlength{\parskip}{2pt}
\setlength{\emergencystretch}{1.5em}
\usepackage{fontspec}
\setmainfont{texgyreheros-regular.otf}[BoldFont=texgyreheros-bold.otf,ItalicFont=texgyreheros-italic.otf,BoldItalicFont=texgyreheros-bolditalic.otf]
\setsansfont{texgyreheros-regular.otf}[BoldFont=texgyreheros-bold.otf]
\usepackage{amsmath,amssymb,booktabs,array,tabularx}
\usepackage{enumitem,xcolor,colortbl,titlesec,fancyhdr,float,caption}
\usepackage[hidelinks]{hyperref}
\definecolor{NYUviolet}{HTML}{57068C}
\definecolor{NYUviolet2}{HTML}{8900E1}
\definecolor{inkgrey}{HTML}{3A3A3A}
\definecolor{rulegrey}{HTML}{B9B9B9}
\definecolor{tablehead}{HTML}{ECE3F4}
\hypersetup{colorlinks=true,linkcolor=NYUviolet,urlcolor=NYUviolet2,citecolor=NYUviolet,
 pdftitle={Assignment 6: GBP-Collateralized FX Forward and Rate Products},pdfauthor={Ziyi Yang}}
\titleformat{\section}{\sffamily\Large\bfseries\color{NYUviolet}}{\thesection}{0.7em}{}[\vspace{2pt}{\color{rulegrey}\titlerule[0.8pt]}]
\titleformat{\subsection}{\sffamily\large\bfseries\color{inkgrey}}{\thesubsection}{0.6em}{}
\titlespacing*{\section}{0pt}{16pt}{8pt}
\titlespacing*{\subsection}{0pt}{10pt}{5pt}
\pagestyle{fancy}
\fancyhf{}
\setlength{\headheight}{15pt}
\renewcommand{\headrulewidth}{0.6pt}
\fancyhead[L]{\footnotesize\sffamily\color{inkgrey}FX Collateralization and Rate Products}
\fancyhead[R]{\footnotesize\sffamily\color{inkgrey}FRE-GT-9743 -- Assignment 6}
\fancyfoot[C]{\footnotesize\sffamily\color{inkgrey}\thepage}
\setlist[itemize]{leftmargin=1.4em,itemsep=2pt,topsep=3pt,label=\textcolor{NYUviolet}{\textbullet}}
\newcolumntype{L}[1]{>{\raggedright\arraybackslash}p{#1}}
\newcolumntype{Y}{>{\raggedright\arraybackslash}X}
\renewcommand{\arraystretch}{1.12}
\newcommand{\keynote}[1]{\par\medskip\noindent\textbf{\sffamily\color{NYUviolet}Key result.}\enspace #1\par\medskip}
\newcommand{\focus}[1]{\noindent{\small\itshape\color{inkgrey}\textbf{Question focus.} #1}\par\medskip}
\numberwithin{equation}{section}
\captionsetup{font=small,labelfont={bf,sf,color=NYUviolet},labelsep=period}

\begin{document}
\setlength{\abovedisplayskip}{7pt}
\setlength{\belowdisplayskip}{7pt}
\setlength{\jot}{2pt}
\hypersetup{pageanchor=false}

\begin{titlepage}
\thispagestyle{empty}
\centering
\vspace*{0.55in}
{\color{NYUviolet}\rule{\textwidth}{2.2pt}}\\[10pt]
{\sffamily\fontsize{16}{19}\selectfont\color{inkgrey}Fixed Income Analysis Report}\\[14pt]
{\sffamily\bfseries\fontsize{25}{30}\selectfont\color{NYUviolet}
GBP-Collateralized FX Forward\\[3pt]and Interest-Rate Products}\\[14pt]
{\sffamily\fontsize{15}{18}\selectfont\color{inkgrey}
Currency-Chain Discounting, Fair Delivery Price,\\[3pt]
RFR Swaps and Overnight Basis Swaps}\\[10pt]
{\color{NYUviolet}\rule{\textwidth}{2.2pt}}\\[26pt]
\begin{minipage}{0.88\textwidth}
\centering\itshape\small\color{inkgrey}
This report derives the value and fair strike of an EUR/USD forward
collateralized in GBP, establishes when the supplied discount factors are
sufficient through a USD-centered currency chain, and documents the completed
implementation of RFR swaps, overnight-index basis swaps, convention-driven
product factories and composite valuation engines.
\end{minipage}\\[34pt]
\begin{tabular}{rl}
\sffamily\bfseries Author & Ziyi Yang\\[3pt]
\sffamily\bfseries Course & FRE-GT-9743\\[3pt]
\sffamily\bfseries Source problem & Jianing (Jay) Yao, Assignment 6\\[3pt]
 & October 19, 2026
\end{tabular}
\vfill
{\sffamily\small\color{inkgrey}Assignment 6 \quad\textbullet\quad FX and Linear Rate Products}\\[2pt]
{\color{NYUviolet}\rule{\textwidth}{1.2pt}}
\end{titlepage}
\hypersetup{pageanchor=true}

\begin{center}{\sffamily\large\bfseries\color{NYUviolet}Abstract}\end{center}
\begin{quote}\small\itshape
An FX forward is the difference between two deterministic currency claims.
Under GBP collateralization, the EUR claim and the USD claim must both be
valued with GBP-collateralized effective discount factors. Their ratio gives
the fair forward strike. Under the course assumption that an FX forward is
independent of collateral currency, a USD-centered currency chain constructs
both required curves from the supplied SONIA, SOFR and USD-collateralized
cross-currency factors. The programming component builds and values fixed-
versus-overnight swaps and SOFR-versus-Fed-Funds basis swaps.
\end{quote}
\vspace{4pt}
{\small\noindent\textbf{\sffamily Keywords:}\;FX forward \textbullet\; Foreign collateral
\textbullet\; Currency chain \textbullet\; RFR swap \textbullet\; Basis swap}
\vspace{8pt}{\color{rulegrey}\hrule}
\setcounter{tocdepth}{1}
\renewcommand{\contentsname}{\sffamily\large\bfseries\color{NYUviolet}Contents}
{\small\tableofcontents}
{\color{rulegrey}\hrule}
\clearpage

\section{Setup and Notation}
Let $S_t^{\$/\mathrm{EUR}}$ be the spot exchange rate in USD per EUR. A long
forward with EUR notional $N$ receives $N$ EUR and pays $NK$ USD at maturity
$T$. Let $P^{x,c}(t,T)$ denote the time-$t$ value, in currency $x$, of one unit
of currency $x$ paid at $T$ when the collateral currency is $c$.

The position consists of two deterministic claims:
\begin{itemize}
\item receive one EUR at $T$, whose USD value today is
$S_t^{\$/\mathrm{EUR}}P^{\mathrm{EUR},\mathrm{GBP}}(t,T)$;
\item pay $K$ USD at $T$, whose USD value today is
$KP^{\$,\mathrm{GBP}}(t,T)$.
\end{itemize}

\section{Q1: Valuation Formula}
\focus{Value the EUR/USD forward in USD when the CSA collateral currency is GBP.}

Combining the two claims gives
\begin{equation}
\boxed{
V_t^{\$}=N\left[
S_t^{\$/\mathrm{EUR}}P^{\mathrm{EUR},\mathrm{GBP}}(t,T)
-KP^{\$,\mathrm{GBP}}(t,T)
\right].}
\end{equation}
Define the GBP-collateralized forward rate by
\begin{equation}
F_t^{\$/\mathrm{EUR};\mathrm{GBP}}(T)
=S_t^{\$/\mathrm{EUR}}
\frac{P^{\mathrm{EUR},\mathrm{GBP}}(t,T)}
{P^{\$,\mathrm{GBP}}(t,T)}.
\end{equation}
Then the mark-to-market takes the familiar form
\begin{equation}
\boxed{V_t^{\$}=NP^{\$,\mathrm{GBP}}(t,T)
\left[F_t^{\$/\mathrm{EUR};\mathrm{GBP}}(T)-K\right].}
\end{equation}
The collateral currency appears in both discount factors. Replacing either
one with an ordinary domestic curve would mix inconsistent collateral regimes.

\section{Q2: Fair Forward Strike}
At inception, the fair strike sets the forward value to zero. Hence
\begin{equation}
S_t^{\$/\mathrm{EUR}}P^{\mathrm{EUR},\mathrm{GBP}}(t,T)
=K_{\mathrm{fair}}P^{\$,\mathrm{GBP}}(t,T),
\end{equation}
and therefore
\begin{equation}
\boxed{
K_{\mathrm{fair}}
=S_t^{\$/\mathrm{EUR}}
\frac{P^{\mathrm{EUR},\mathrm{GBP}}(t,T)}
{P^{\$,\mathrm{GBP}}(t,T)}.}
\end{equation}
The numerator discounts the received base-currency claim and the denominator
discounts the paid quote-currency claim under the same GBP collateral rule.

\section{Q3: Sufficiency of the Supplied Curves}
The curve set provides
\begin{equation}
P^{\mathrm{GBP},\mathrm{GBP}}=df_{\mathrm{SONIA}},\quad
P^{\mathrm{EUR},\mathrm{EUR}}=df_{\mathrm{EuroSTR}},\quad
P^{\$,\$}=df_{\mathrm{SOFR}},
\end{equation}
together with $P^{\mathrm{GBP},\$}=df[\mathrm{GBP},\$]$ and
$P^{\mathrm{EUR},\$}=df[\mathrm{EUR},\$]$.

Under the course currency-chain assumption that the same-maturity FX forward
is independent of the collateral currency, a USD-centered chain gives
\begin{align}
P^{\mathrm{EUR},\mathrm{GBP}}
&=\boxed{\frac{P^{\mathrm{EUR},\$}}{P^{\mathrm{GBP},\$}}
P^{\mathrm{GBP},\mathrm{GBP}}},\\
P^{\$,\mathrm{GBP}}
&=\boxed{\frac{P^{\$,\$}}{P^{\mathrm{GBP},\$}}
P^{\mathrm{GBP},\mathrm{GBP}}}.
\end{align}
Both GBP-collateralized factors required by the valuation formula can therefore
be constructed from the supplied set. In the fair-strike ratio, the GBP terms
and the GBP/USD cross-currency factor cancel:
\begin{equation}
\boxed{
K_{\mathrm{fair}}
=S_t^{\$/\mathrm{EUR}}
\frac{P^{\mathrm{EUR},\$}}{P^{\$,\$}}.}
\end{equation}

\keynote{The supplied curves are sufficient under the assignment's currency-chain assumption. The EUR domestic $df_{\mathrm{EuroSTR}}$ is not required for this particular forward once $P^{\mathrm{EUR},\$}$ is available.}

The conclusion is conditional. If forward prices can depend materially on the
collateral currency because of CSA optionality, liquidity, transaction costs or
convexity effects, the chain identity is not automatic. Then direct
$P^{\mathrm{EUR},\mathrm{GBP}}$ and $P^{\$,\mathrm{GBP}}$ curves, or equivalent
GBP-collateralized cross-currency basis data, are needed.

\section{Completed Product Library}
\subsection{RFR swap}
The completed product constructor builds a fixed stream and a compounded
overnight stream with their contractual accrual periods and payment
conventions. For receive-fixed,
\begin{equation}
PV_{\mathrm{swap}}=PV_{\mathrm{fixed}}-PV_{\mathrm{floating}}.
\end{equation}
The floating coupons use the SOFR projection curve, while both legs are
discounted with the selected USD funding curve. The same signs are applied to
cash paid or received on the value date.

\subsection{Overnight-index basis swap}
Leg 1 contains compounded SOFR plus the quoted spread; leg 2 contains
compounded Fed Funds. Each index uses its own projection curve. The
implementation also exposes the leg-1 index and spread as separate streams,
so the valuation engine combines
\begin{equation}
PV_{\mathrm{basis}}
=PV_{\mathrm{leg1,index}}+PV_{\mathrm{leg1,spread}}
-PV_{\mathrm{leg2,index}}
\end{equation}
when leg 1 is received.

\subsection{Convention-driven factory}
The two completed builders map stored conventions into product constructors.
For swaps, \texttt{3Mx1Y} means a start three months after the value date and
a one-year life from that start; \texttt{1Y} means spot-starting for one year.
The quote is the fixed rate for an RFR swap and the spread over leg 1 for a
basis swap. The notebook verifies that factory-built products serialize
identically to directly built products.

\begin{table}[H]
\centering
\caption{Verified Assignment 6 product values}
\small
\begin{tabularx}{\textwidth}{@{}L{4.1cm}L{3cm}Y@{}}
\toprule\rowcolor{tablehead}
\textbf{Product}&\textbf{PV}&\textbf{Verification}\\\midrule
FRA & $87.090294$ & Matches the analytical forward-rate formula\\
Receive-fixed RFR swap & $-2{,}629.679575$ & Fixed-leg PV minus floating-leg PV\\
Receive-leg-1 basis swap & $8{,}446.472471$ & SOFR-plus-spread PV minus Fed-Funds PV\\
\bottomrule
\end{tabularx}
\end{table}

The executed notebook additionally checks a known FRA fixing, cashflow reports,
factory equality and the public valuation API. All numerical assertions pass.

\section{Conclusion}
Foreign-collateral FX pricing requires one effective discount factor for each
settlement currency under the common CSA currency. The fair strike is their
ratio times spot. Under the stated currency-chain assumption, the supplied
USD-centered curves are sufficient to reconstruct both GBP-collateralized
factors. The completed code applies the same decomposition principle to linear
rate products: construct contractual cashflow streams, project each floating
leg with its index curve, discount with the funding curve and combine the legs
with the correct receive/pay signs.

\end{document}
